\section{Reduzierfunktion}
Dieser Abschnitt widmet sich der Reduzierfunktion. Zunächst wird der Funktionshintergrund und die Funktionsweise dieser Funktion beschrieben. Im Anschluss erfolgt eine detaillierte Analyse der spezifischen Anforderungen sowie der bestehenden Softwareimplementierung (C-Code). Auf Grundlage der gewonnenen Erkenntnisse wird die Schnittstelle für den Playground erweitert und die Funktion in Matlab modelliert.

\subsection{Funktionsbeschreibung}
Die Reduzierfunktion umfasst die Begrenzung der maximalen Leistung, des maximalen Motorstroms und der maximalen Drehzahl des Geräts. Sie erfüllt primär drei funktionale Ziele:
\begin{enumerate}
\item Schonung des Akkus und Optimierung der Restkapazität: Durch die Reduktion der Leistung und des Stroms bei Erreichen einer definierten Akkuspannung wird der Akku entlastet. Dies ermöglicht eine effektivere Nutzung der verbleibenden Kapazität und trägt zur Schonung des Akkus bei.
\item Nutzerfeedback: Die Reduzierung der maximalen Drehzahl dient als Hinweis für den Nutzer auf den sinkenden Ladezustand des Akkus und signalisiert, dass eine baldige Aufladung erforderlich ist.
\item Thermische Entlastung: Eine Begrenzung von Leistung und Strom kann zur thermischen Stabilisierung des Geräts beitragen und schützt es vor Überhitzung.
\end{enumerate}
Die Aktivierung der Reduzierfunktion erfolgt durch drei verschiedene Triggerbedingungen:
\begin{enumerate}
\item Unterschreiten der Akkuspannung: Sobald die Akkuspannung unter einen vordefinierten Schwellenwert fällt, wird die Reduzierfunktion aktiviert.
\item Erhöhte Elektroniktemperatur: Ein Anstieg der Elektroniktemperatur über einen festgelegten Grenzwert führt ebenfalls zur Aktivierung der Funktion.
\item Erhöhte Motortemperatur: Auch eine Überschreitung des definierten Temperaturgrenzwerts für den Motor löst die Reduzierfunktion aus.
\end{enumerate}
Die Deaktivierung der Reduzierfunktion erfolgt, wenn die jeweiligen Bedingungen wieder in den Normalbereich zurückkehren.\\
Im Folgenden werden die Auswirkungen der Reduzierfunktion anhand der Leistungskennlinie erläutert. Eine Reduktion des Motorstroms führt zu einer Abflachung des ansteigenden Abschnitts der Kennlinie. Die prozentuale Verringerung der Leistung bewirkt eine Absenkung des Plateaus, während die Reduktion der Drehzahl zu einer Verkürzung der Kennlinie führt \cite{SW_ZUSE}. Diese Zusammenhänge sind in \autoref{png_auswirkungen_red} schematisch dargestellt.
\begin{figure}[H]
	\centering
	\includegraphics[width=0.8\textwidth]{./Bilder/Funktionsbeschreibung_Red.png}
	\caption{Auswirkung der Reduzierfunktion auf die Leistungskennlinie \cite{SW_ZUSE}}
	\label{png_auswirkungen_red}
\end{figure}

\subsection{Analyse des Lastenhefts und des C-Codes}
\label{subsec_red_analyse}
In diesem Abschnitt werden die spezifischen Anforderungen aus dem Lastenheft bei STIHL für die Reduzierfunktion aufgeführt. Diese Anforderungen bilden die Grundlage für die Implementierung des entsprechenden C-Codes in der Firmware. Durch eine ergänzende Analyse des C-Codes kann die Modellierung der Funktion vereinfacht und gleichzeitig sichergestellt werden, dass deren Funktionalität vollständig abgebildet wird. \\
Die nachfolgenden Anforderung sind wörtlich aus dem Master-Lastenheft für Akkugeräte entnommen \cite{Lastenheft}:
\begin{itemize}
\item Wenn die Akkuspannung den Wert \textit{CpsRdzSpgDrehzahlReduzierungEinZelle} für die Zeit \textit{CpsRdzZeiDrehzahlReduzierungEin} unterschreitet, soll die Reduzierung aktiviert werden.
\item Wenn die Reduzierung aufgrund der Akkuspannung aktiv ist und die Akkuspannung den Wert \textit{CpsRdzSpgReduzierungAusZelle} für die Zeit \textit{CpsRdzZeitReduzierungAus} überschreitet und der Motor steht, soll das Gerät die Reduzierung aufgrund der Akkuspannung deaktivieren.
\item Überschreitet die gemessene Elektroniktemperatur den Wert \textit{CpsRdzTmcElektronikEin} oder die berechnete Motortemperatur den Wert \textit{CpsRdzTmcMotorEin}, so wird die Reduzierungsfunktion aktiviert.
\item Sinkt die Elektroniktemperatur unter \textit{CpsRdzTmcElektronikAus} und die Motortemperatur unter \textit{CpsRdzTmcMotorAus} und steht der Motor, so wird die Reduktion von Leistung, Strom und Drehzahl aufgrund dieser Funktion beendet.
\item Wenn die Reduzierung aufgrund Akkuspannung oder Temperatur aktiviert ist, soll die maximale Leistung und die maximale Drehzahl prozentual (\textit{CpsRdzDutDeltaLeistungReduzierung, CpsRdzDutDrehzahlDeltaReduzierung}) zur aktuellen Stufe/Modus sowie der maximale Motorstrom auf einen fixen Wert \textit{CpsRdzStrMotorMaxReduzierung} reduziert werden.
\end{itemize}
Der C-Code in der Firmware basiert auf den genannten Anforderungen. Eine zusätzliche Untersuchung des C-Codes ermöglicht es, sowohl die Informationen zum Ein- und Ausschalten der Funktion über maschinenabhängige Optionen als auch die zusätzlichen Eingangssignale, die für eine optimale Funktionalität der Funktion erforderlich sind, zu identifizieren. \\
Aus der Analyse des C-Codes \cite{Firmware_gesamt} ergeben sich folgende Punkte für die Modellierung der Reduzierfunktion:
\begin{itemize}
\item Die Funktion ist nicht über eine maschinenabhängige Option ein- oder ausschaltbar, sondern standardmäßig bei jeder Maschine aktiviert.
\item Die Reduzierfunktion erfordert für ihre Aktivierung sowohl das Validierungsflag für die Temperaturen als auch ein zusätzliches Validierungsflag für die Akkuspannung, welches in der Basissoftware implementiert ist. Die Funktion darf nur aktiviert sein, wenn beide Flags gesetzt sind.
\item Die Parameter zur Reduzierung von Drehzahl, Leistung und Strom variieren je nach den spezifischen Bedingungen für die Aktivierung der Funktion. Daher ergibt sich für die Strombegrenzung bei einer Reduzierung aufgrund der Motortemperatur ein anderer Wert als bei einer Reduzierung aufgrund der Elektroniktemperatur. Ist die Reduzierfunktion durch mehrere Bedingungen gleichzeitig aktiv, wird der jeweils minimale Wert für Drehzahl, Leistung und Strom verwendet.
\item Bei Akkugeräten mit zwei Akkus existiert eine Sonderfunktion für die Reduktion aufgrund der Akkuspannung. Ist diese Reduktion aktiv und verfügt der zweite Akku über eine ausreichend hohe Spannung, wird die Reduktion unmittelbar zurückgesetzt, ohne die Wartezeit gemäß \textit{CpsRdzZeitReduzierungAus} abzuwarten.
\item Für Geräte mit einem HMI besteht die Möglichkeit, über eine LED anzuzeigen, ob eine Reduzierung durch die Motortemperatur oder die Elektroniktemperatur ausgelöst wurde. Hierfür muss ein Status-Flag an die Basissoftware übermittelt werden, das den aktuellen Zustand der Reduzierung anzeigt.
\item Die Funktion wird in einem Intervall von \unit[1]{ms} zyklisch ausgeführt.
\end{itemize}
\subsection{Erweiterung der Schnittstelle für den Playground}
In diesem Abschnitt wird die Erweiterung der Schnittstelle zur korrekten Übergabe der Ein- und Ausgangssignale zwischen der Basissoftware und dem Autocode behandelt. Zu diesem Zweck wird die Datei \textit{Playground.c} angepasst. Die erforderlichen Eingangssignale werden aus der Basissoftware in der Datei \textit{Playground.c} aufgerufen und anschließend an den Playground übergeben. Der entsprechende C-Code wird in dieser Arbeit nicht dargestellt. Stattdessen ist in \autoref{png_io_red} ein schematisches Diagramm für die Reduzierfunktion mit allen Ein- und Ausgängen abgebildet.
\begin{figure}[H]
	\centering
	\includegraphics[width=0.8\textwidth]{./Bilder/IO_Red.pdf}
	\caption{Ein- und Ausgänge der Reduzierfunktion}
	\label{png_io_red}
\end{figure}
Die Reduzierfunktion benötigt als Eingänge, analog zu den zuvor beschriebenen Funktionen, die Motor- und Elektroniktemperatur sowie das zugehörige Validierungsflag. Zusätzlich ist die aktuelle Akkuspannung samt Validierungsflag erforderlich. Weiterhin werden die maximale Leistung und die maximale Drehzahl der aktuellen Betriebsstufe als Eingangsgrößen benötigt, um im Fall einer aktiven Reduzierung die entsprechend reduzierten Werte zu berechnen. Wie in \autoref{subsec_red_analyse} beschrieben, darf die Rücksetzung der Reduzierung, ausschließlich im Stillstand des Motors erfolgen. Dafür stellt die Basissoftware einen Motorstatus zur Verfügung, der angibt, ob sich die Maschine im \textit{running}- oder \textit{off}-Modus befindet.\\
Die Funktion liefert als Ausgänge den maximalen Strom, die berechnete maximale Leistung sowie die maximale Drehzahl. Zusätzlich wird ein Flag für ein mögliches HMI bereitgestellt, das den Status einer aktiven Reduzierung aufgrund von einer Temperaturüberschreitung (Motor oder Elektronik) anzeigt.

\subsection{Modellierung der Funktion in Matlab}
Dieser Abschnitt behandelt die Modellierung der Reduzierfunktion in Matlab. Wie in \autoref{ch_voruntersuchungen} erläutert, wird für jede Teilfunktion ein eigenes Referenced-Subsystem sowie ein individuelles Simulink Data Dictionary erstellt. Im Dictionary werden alle Parameter entsprechend den Anforderungen (\autoref{subsec_red_analyse}) definiert. Dabei erfolgt die Benennung der Parameter gemäß der in \autoref{subsec_namenskonvention} beschriebenen Namenskonvention. Zusätzlich werden zentrale Signale, wie beispielsweise die Ausgänge der Funktion, ebenfalls im Dictionary hinterlegt. Dies ermöglicht den Zugriff auf die Signale über eine A2L-Datei und damit das Messen der Signale auf der realen Maschine über eine XCP-Schnittstelle. \\
Die vollständige Modellierung der Reduzierfunktion ist in \autoref{png_matlab_red_erk_über} und \ref{png_matlab_red_aus} dargestellt.
\begin{figure}[H]
	\centering
	\includegraphics[width=1\textwidth]{./Bilder/Modell_Red_Erk_übersicht.pdf}
	\caption{Matlab-Modell: Reduziererkennung}
	\label{png_matlab_red_erk_über}
\end{figure}
\begin{figure}[H]
	\centering
	\includegraphics[width=1\textwidth]{./Bilder/Modell_Red_Ausw.pdf}
	\caption{Matlab-Modell: Reduzierauswertung}
	\label{png_matlab_red_aus}
\end{figure}
Im oberen Bereich der Darstellung ist das Stateflow-Chart für die Reduziererkennung abgebildet, in dem die einzelnen Flags zur Aktivierung oder Deaktivierung der Reduzierfunktion gesetzt werden. Im unteren Bereich ist die Auswertung dieser Flags modelliert.\\
Auf der linken Seite verfügt jedes Flag über einen eigenen Schalter, der zwischen aktivierter und deaktivierter Reduzierung umschaltet. Ist die Reduzierung deaktiviert, wird der \textit{BASIS\_CURRENT} als maximaler Strom genutzt, während für die prozentuale Reduzierung von Leistung und Drehzahl ein Wert von null Prozent eingesetzt wird. Bei aktivierter Reduzierung werden die Parameter aus dem Lastenheft verwendet. In der Mitte der Darstellung ermitteln Min-/Max-Blöcke den kleinsten maximalen Strom sowie den höchsten Prozentwert für Leistung und Drehzahl. Basierend auf diesen Werten werden die reduzierte maximale Leistung und die reduzierte maximale Drehzahl in Abhängigkeit der Eingangswerte \textit{max\_power} und \textit{max\_speed} berechnet.\\
Die berechneten Größen werden zu einem Bus-Objekt zusammengeführt und als Ausgang des Subsystems bereitgestellt. Zusätzlich wird aus den beiden Flags, die für die Reduzierung aufgrund der Temperaturen zuständig sind, durch eine logische Oder-Verknüpfung das Flag für das HMI gesetzt.

Innerhalb der Reduziererkennung ist die Funktion nur dann aktiv, wenn sowohl das Flag für die Validierung der Temperaturen als auch das Flag für die Validierung der Akkuspannung gesetzt sind. Sind diese Bedingungen erfüllt, werden drei Subcharts parallel ausgeführt, die jeweils für die Erkennung der Reduzierung aufgrund der Akkuspannung, der Motortemperatur und der Elektroniktemperatur verantwortlich sind. Diese Struktur ist in \autoref{png_matlab_red_erk} dargestellt.
\begin{figure}[H]
	\centering
	\includegraphics[width=1\textwidth]{./Bilder/Modell_Red_Erk.pdf}
	\caption{Matlab-Modell: Stateflow-Chart für die Reduziererkennung}
	\label{png_matlab_red_erk}
\end{figure}
Der Inhalt der einzelnen Subcharts wird im Rahmen dieser Arbeit nicht näher erläutert.

%\begin{figure}[H]
%	\centering
%	\includegraphics[width=1\textwidth]{./Bilder/Modell_Red_Erk_spngAkku.pdf}
%	\caption{}
%	\label{png_matlab_red_erk_spng}
%\end{figure}
%
%\begin{figure}[H]
%	\centering
%	\includegraphics[width=1\textwidth]{./Bilder/Modell_Red_Erk_tmpElk.pdf}
%	\caption{}
%	\label{png_matlab_red_erk_tmpelk}
%\end{figure}
%
%\begin{figure}[H]
%	\centering
%	\includegraphics[width=1\textwidth]{./Bilder/Modell_Red_Erk_tmpMot.pdf}
%	\caption{}
%	\label{png_matlab_red_erk_tmpmot}
%\end{figure}